\documentclass[11pt]{article}

\usepackage[letterpaper,margin=0.82in,headheight=15pt]{geometry}
\usepackage[T1]{fontenc}
\usepackage{lmodern}
\usepackage{microtype}
\usepackage{amsmath,amssymb,amsthm,bm,mathtools}
\usepackage{stmaryrd}
\usepackage{booktabs,longtable,tabularx,array,multirow}
\usepackage{enumitem}
\usepackage{xcolor}
\usepackage{fancyhdr}
\usepackage{graphicx}
\usepackage{tikz-cd}
\usepackage{hyperref}
\usepackage[nameinlink,noabbrev]{cleveref}

\definecolor{navy}{HTML}{17324D}
\definecolor{teal}{HTML}{147D92}
\definecolor{orange}{HTML}{D47A27}
\definecolor{softblue}{HTML}{EEF4F7}
\definecolor{softgray}{HTML}{F3F5F6}
\definecolor{darkgray}{HTML}{46545E}
\definecolor{softorange}{HTML}{FFF4E8}
\definecolor{softgreen}{HTML}{EDF7F0}
\definecolor{softred}{HTML}{FCEEEE}

\hypersetup{
  colorlinks=true,
  linkcolor=navy,
  citecolor=teal,
  urlcolor=teal,
  pdftitle={Volume IX: Dynamical Gluon Fock Sectors, Sector-Dependent Renormalization, and Gauge-Consistent Microscopic Wilson Dynamics},
  pdfauthor={DeuteronWigner project}
}

\pagestyle{fancy}
\fancyhf{}
\fancyhead[L]{\small\color{darkgray}DeuteronWigner formalism - Volume IX}
\fancyhead[R]{\small\color{darkgray}Dynamical gluon sectors and gauge closure}
\fancyfoot[L]{\small\color{darkgray}30 July 2026}
\fancyfoot[R]{\small\color{darkgray}\thepage}
\renewcommand{\headrulewidth}{0.4pt}

\setlist[itemize]{leftmargin=1.5em,itemsep=2pt,topsep=3pt}
\setlist[enumerate]{leftmargin=1.7em,itemsep=2pt,topsep=3pt}
\setlength{\parindent}{0pt}
\setlength{\parskip}{5pt}
\setlength{\emergencystretch}{2em}
\renewcommand{\arraystretch}{1.16}
\allowdisplaybreaks

\newtheorem{definition}{Definition}[section]
\newtheorem{axiom}[definition]{Axiom}
\newtheorem{principle}[definition]{Design principle}
\newtheorem{requirement}[definition]{Requirement}
\newtheorem{proposition}[definition]{Proposition}
\newtheorem{theorem}[definition]{Theorem}
\newtheorem{lemma}[definition]{Lemma}
\newtheorem{corollary}[definition]{Corollary}
\newtheorem{remark}[definition]{Remark}
\newtheorem{decision}[definition]{Model decision}

\crefname{definition}{definition}{definitions}
\crefname{axiom}{axiom}{axioms}
\crefname{principle}{design principle}{design principles}
\crefname{requirement}{requirement}{requirements}
\crefname{proposition}{proposition}{propositions}
\crefname{theorem}{theorem}{theorems}
\crefname{lemma}{lemma}{lemmas}
\crefname{corollary}{corollary}{corollaries}
\crefname{remark}{remark}{remarks}
\crefname{decision}{model decision}{model decisions}

\newcommand{\kT}{\bm{k}_{T}}
\newcommand{\pT}{\bm{p}_{T}}
\newcommand{\qT}{\bm{q}_{T}}
\newcommand{\DT}{\bm{\Delta}_{T}}
\newcommand{\bD}{\bm{b}_{\Delta}}
\newcommand{\bM}{\bm{b}_{\mathrm{TMD}}}
\newcommand{\dd}{\mathrm{d}}
\newcommand{\Tr}{\operatorname{Tr}}
\newcommand{\tr}{\operatorname{tr}}
\newcommand{\diag}{\operatorname{diag}}
\newcommand{\rank}{\operatorname{rank}}
\newcommand{\Span}{\operatorname{span}}
\newcommand{\Ker}{\operatorname{Ker}}
\newcommand{\Ran}{\operatorname{Ran}}
\newcommand{\Cov}{\operatorname{Cov}}
\newcommand{\Var}{\operatorname{Var}}
\newcommand{\Id}{\operatorname{Id}}
\newcommand{\Hom}{\operatorname{Hom}}
\newcommand{\End}{\operatorname{End}}
\newcommand{\Hil}{\mathcal H}
\newcommand{\LF}{\ensuremath{\mathrm{LF}}}
\newcommand{\BLFQ}{\ensuremath{\mathrm{BLFQ}}}
\newcommand{\TMD}{\ensuremath{\mathrm{TMD}}}
\newcommand{\GTMD}{\ensuremath{\mathrm{GTMD}}}
\newcommand{\GPD}{\ensuremath{\mathrm{GPD}}}
\newcommand{\QCD}{\ensuremath{\mathrm{QCD}}}
\newcommand{\EFT}{\ensuremath{\mathrm{EFT}}}
\newcommand{\TTN}{\ensuremath{\mathrm{TTN}}}
\newcommand{\TN}{\ensuremath{\mathrm{TN}}}
\newcommand{\cO}{\mathcal O}
\newcommand{\cM}{\mathcal M}
\newcommand{\cR}{\mathcal R}
\newcommand{\cS}{\mathcal S}
\newcommand{\cT}{\mathcal T}
\newcommand{\cV}{\mathcal V}
\newcommand{\cW}{\mathcal W}
\newcommand{\cE}{\mathcal E}
\newcommand{\cG}{\mathcal G}
\newcommand{\cK}{\mathcal K}
\newcommand{\cQ}{\mathcal Q}
\newcommand{\cP}{\mathcal P}
\newcommand{\cC}{\mathcal C}
\newcommand{\cZ}{\mathcal Z}
\newcommand{\cD}{\mathcal D}
\newcommand{\cN}{\mathcal N}
\newcommand{\cU}{\mathcal U}
\newcommand{\cF}{\mathcal F}
\newcommand{\cA}{\mathcal A}
\newcommand{\cB}{\mathcal B}
\newcommand{\cL}{\mathcal L}
\newcommand{\cJ}{\mathcal J}
\newcommand{\cI}{\mathcal I}
\newcommand{\eps}{\varepsilon}
\newcommand{\abs}[1]{\left|#1\right|}
\newcommand{\norm}[1]{\left\lVert#1\right\rVert}
\newcommand{\inner}[2]{\left\langle #1\middle|#2\right\rangle}
\newcommand{\avg}[1]{\left\langle #1\right\rangle}
\newcommand{\phys}{\mathrm{phys}}
\newcommand{\eff}{\mathrm{eff}}
\newcommand{\bare}{\mathrm{bare}}
\newcommand{\ind}{\mathrm{ind}}
\newcommand{\ct}{\mathrm{ct}}
\newcommand{\reg}{\mathrm{reg}}
\newcommand{\trunc}{\mathrm{trunc}}
\newcommand{\ren}{\mathrm{ren}}
\newcommand{\num}{\mathrm{num}}
\newcommand{\UV}{\mathrm{UV}}
\newcommand{\IR}{\mathrm{IR}}
\newcommand{\COM}{\mathrm{CM}}
\newcommand{\Lz}{L_z}
\newcommand{\Jz}{J^z}
\newcommand{\Pthree}{P_3}
\newcommand{\Pfour}{P_{4g}}
\newcommand{\Hthree}{\Hil_{qqq}}
\newcommand{\Hfour}{\Hil_{qqqg}}
\newcommand{\Mtwo}{\cM^2}
\newcommand{\WTI}{\mathrm{WTI}}
\newcommand{\STI}{\mathrm{STI}}
\newcommand{\PV}{\operatorname{PV}}
\newcommand{\AssumptionBundle}{\texttt{AssumptionBundle}}
\newcommand{\PredictionPlan}{\texttt{PredictionPlan}}
\newcommand{\StateBundle}{\texttt{MicroscopicStateBundle}}
\newcommand{\WardManifest}{\texttt{WardClosureManifest}}
\newcommand{\WilsonHandoff}{\texttt{MicroscopicWilsonHandoff}}
\newcommand{\RenTrajectory}{\texttt{SectorRenormalizationTrajectory}}
\newcommand{\TermRegistry}{\texttt{HamiltonianTermRegistry}}
\newcommand{\ProvComplex}{\texttt{Provenance2Complex}}
\newcommand{\Amp}{\mathbf{Amp}}
\newcommand{\Match}{\mathbf{Match}}
\newcommand{\Red}{\mathbf{Red}}
\newcommand{\Proc}{\mathbf{Proc}}
\newcommand{\Infer}{\mathbf{Infer}}

\newcolumntype{Y}{>{\raggedright\arraybackslash}X}
\newcolumntype{L}[1]{>{\raggedright\arraybackslash}p{#1}}

\title{\vspace{-1.0cm}
{\color{navy}\bfseries Volume IX: Dynamical Gluon Fock Sectors,\\[2mm]
Sector-Dependent Renormalization, and Gauge-Consistent\\[1mm]
Microscopic Wilson Dynamics}\\[4mm]
\large The coupled \(qqq\oplus qqqg\) nucleon, instantaneous partners, Ward closure,
gluon/OAM state exports, and the microscopic bridge to GTMD and Wilson operators}
\author{DeuteronWigner project}
\date{30 July 2026}

\begin{document}
\maketitle

\begin{center}
\begin{minipage}{0.94\textwidth}
\colorbox{softblue}{\parbox{0.96\textwidth}{
\textbf{Status and purpose.}
Volumes~0--VIII established the algebraic and geometric architecture, regulated
light-front Hilbert spaces, common GTMD operators, dynamical one-gluon Wilson
pilots, matched spin-1 nuclear structure, QCD matching and evolution contracts,
shared inference, a concrete RSA--BLFQ Hamiltonian route, and a
symmetry-adapted tensor-network prediction compiler.  C7/H0 validated the
microscopic basis and Hamiltonian-term spine.  C8/H1 then produced the first
interacting valence eigenstate, resolution-dependent renormalization flow,
Hamiltonian-consistent vector current, exact/Krylov solvers, and a genuine
variational tree tensor network.  C9/H2 is the active implementation package
for the first coupled \(qqq\oplus qqqg\) eigenproblem.  The present volume is
the normative specification for that transition.  Its central requirement is
that the propagating gluon sector, instantaneous light-front interactions,
sector-dependent counterterms, current attachments, field-strength operators,
and Wilson-line couplings be treated as one gauge-consistent renormalized
object.  An explicit \(qqqg\) sector may not be combined with the valence
operator that already represents its eliminated effects unless a matching and
subtraction cell is supplied.  The resulting state remains a controlled
low-resolution Hamiltonian-EFT state; it is not yet a continuum-QCD nucleon or
a soft-subtracted physical TMD.
}}
\end{minipage}
\end{center}

\begin{abstract}
The first dynamical gluon sector changes the model class more deeply than the
addition of another basis block.  It introduces a gauge-field degree of
freedom into the same interacting state that carries quark flavor, helicity,
orbital angular momentum, color multiplicity, current response, and the
microscopic ingredients of Wilson-line rescattering.  We formulate the
coupled \(qqq\oplus qqqg\) nucleon as a resolution-indexed block
light-front mass operator with two independent \(qqqg\) color-singlet
multiplicities, exact quark antisymmetry, propagating quark--gluon emission and
absorption, instantaneous-fermion and instantaneous-gluon partners,
sector-dependent mass and vertex counterterms, induced confinement, and an
explicit truncation-discrepancy basis.  Gauge consistency is tested first by
an Abelianized finite-basis Ward--Takahashi identity whose propagating,
instantaneous, current, wave-function, and counterterm contributions share one
Hamiltonian identity.  Passing this test does not establish the full
non-Abelian Slavnov--Taylor identities; the corresponding ghost, three-gluon,
higher-Fock, and regulator obligations remain explicit.  The coupled state is
represented by a symmetry-adapted direct-sum tensor network with a retained
Fock-sector edge and outer color multiplicities.  Exact diagonalization,
matrix-free Krylov methods, and variational TTN optimization define a common
solver hierarchy, with convergence tested not only in energy but also in gluon
probability, momentum, helicity, orbital angular momentum, currents, and Ward
residuals.  We derive the finite-resolution quark/gluon momentum and spin
ledgers, the field-strength operator exports needed by gluon GTMDs, the typed
handoff to the validated C5/C6 Wilson engine, and the sector-completeness
conditions for higher Wilson orders.  A provenance two-complex relates the
explicit \(qqqg\) theory to its Feshbach-induced valence reduction and forbids
double counting.  The volume concludes with formal requirements, benchmarks,
negative tests, software schemas, and the interface to the light-sea/chiral
H3 program.
\end{abstract}

\tableofcontents

\section{Scientific purpose and boundary}

\subsection{Why the dynamical gluon sector is decisive}

C8/H1 supplies an interacting valence state
\begin{equation}
 |N,\Lambda\rangle_{\mathrm{H1}}
 \in \Hthree,
\end{equation}
with a resolution-dependent effective Hamiltonian, current, state-tracking
record, and tensor-network representation.  That state can test valence mass,
charge, spin, orbital, and numerical convergence, but it cannot generate a
gluon distribution, a genuine quark--gluon correlator, or a Wilson phase using
the gauge field of the same eigenstate.  The first required extension is
\begin{equation}
 \Hil_{\mathrm{H2},r}
 =
 \Hthree^{(r)}\oplus\Hfour^{(r)},
 \label{eq:h2-space}
\end{equation}
with
\begin{equation}
 |N,\Lambda;r\rangle
 =
 |\Psi_{3}^{\Lambda;r}\rangle
 +
 |\Psi_{4g}^{\Lambda;r}\rangle .
 \label{eq:h2-state}
\end{equation}
This extension makes the following objects dynamical rather than assigned:
\begin{itemize}
 \item the probability and momentum carried by one resolved gluon;
 \item the two independent color-singlet embeddings of \(qqqg\);
 \item the quark--gluon vertex and its Hermitian conjugate;
 \item the OAM and helicity interferences involving the gluon sector;
 \item self-energy and current corrections generated by the same sector;
 \item the state amplitudes consumed by gluon GTMD and Wilson-line operators.
\end{itemize}
Recent BLFQ calculations demonstrate that nucleon eigenstates containing
\(qqq\), \(qqqg\), and light-pair sectors can support quark and gluon density,
helicity, transversity, and imaging calculations
\cite{Mondal2025,Vary2025}.  Those results establish feasibility, not the
identity of the present Hamiltonian or its renormalization trajectory.

\begin{decision}[Primary H2 realization]
The first dynamical-gluon calculation is the RSA--BLFQ Hamiltonian-EFT route of
Volume~VII, enlarged from \(qqq\) to \(qqq\oplus qqqg\) and executed through
the assumption compiler and symmetry-adapted tensor networks of Volume~VIII.
Published BLFQ states are comparison and validation data; they are not silently
substituted for the project state.
\end{decision}

\subsection{Inherited contracts}

H2 inherits the following immutable contracts.
\begin{enumerate}
 \item The light-front convention is
 \begin{equation}
  v^{\pm}=\frac{v^0\pm v^3}{\sqrt2},
  \qquad
  \cM^2=2P^+P^- -\bm P_T^2.
 \end{equation}
 \item The regulator, basis, Fock content, counterterms, and renormalization
 conditions are part of the Hamiltonian identity.
 \item The two \(qqqg\) singlet multiplicity channels remain distinct.
 \item Identical-quark antisymmetry is imposed before matrix assembly.
 \item Hamiltonians, currents, GTMD operators, and Wilson couplings transform
 together under any sector elimination or similarity transformation.
 \item The C5/C6 Wilson engine may consume a microscopic H2 state only through
 a typed handoff; it may not infer cut support or soft completion from state
 existence.
 \item A regulated Hamiltonian overlap is not a physical TMD until LF-to-QCD,
 UV, rapidity, soft, and scheme matching are complete.
\end{enumerate}

\subsection{Explicit nonclaims}

This volume does not claim:
\begin{itemize}
 \item a continuum-QCD solution of the nucleon;
 \item a complete non-Abelian Slavnov--Taylor closure at finite Fock order;
 \item a complete zero-mode solution;
 \item physical gluon PDFs, GTMDs, or TMDs without matching;
 \item physical absorptive cuts from a discrete bound-state spectrum alone;
 \item Wilson orders beyond the Fock sectors that can support them;
 \item a sea-quark or antiquark prediction before the H3 pair sectors;
 \item a deuteron prediction, process factorization theorem, or global fit.
\end{itemize}

\begin{principle}[No gauge closure by selective omission]
A finite-basis Ward test may be declared successful only if every propagating,
instantaneous, current, counterterm, endpoint, and regulator contribution
required at the stated order is present.  Omitting a difficult term and fitting
the remaining current to zero defect is not gauge closure.
\end{principle}

\section{The coupled \texorpdfstring{\(qqq\oplus qqqg\)}{qqq+qqqg} Hilbert space}

\subsection{Sector decomposition and projectors}

Let
\begin{equation}
 \Pthree+\Pfour=\Id_{\mathrm{H2}},
 \qquad
 \Pthree\Pfour=0,
 \qquad
 \Pthree^2=\Pthree,
 \quad
 \Pfour^2=\Pfour .
\end{equation}
The normalized state obeys
\begin{equation}
 P_3^{(r)}+P_{4g}^{(r)}=1,
 \qquad
 P_s^{(r)}
 =
 \langle\Psi_N^{(r)}|P_s|\Psi_N^{(r)}\rangle .
 \label{eq:sector-prob}
\end{equation}
The individual probabilities depend on regulator and resolution.  Their
physical role is diagnostic; matched observables and total quantum-number
ledgers are the convergence targets.

The one-gluon sector basis states have the schematic form
\begin{equation}
 |\alpha_{4g}\rangle
 =
 \cA_{123}
 \left|
 \{k_i,n_i,m_i,\lambda_i,f_i,c_i\}_{i=1}^{3};
 k_g,n_g,m_g,\lambda_g,a;
 \rho,\Jz,I,r
 \right\rangle,
 \label{eq:qqqg-basis}
\end{equation}
where \(\cA_{123}\) is the exact quark antisymmetrizer and \(\rho\) labels the
outer multiplicity of the three-quark color octet that couples with the
adjoint gluon to a singlet.

\subsection{Color singlets and outer multiplicity}

The tensor product contains
\begin{equation}
 3\otimes3\otimes3
 =
 1\oplus 8_{\rho}\oplus 8_{\lambda}\oplus10,
\end{equation}
and therefore
\begin{equation}
 (8_{\rho}\oplus8_{\lambda})\otimes8
 \supset
 1_{\rho}\oplus1_{\lambda}.
 \label{eq:two-singlets}
\end{equation}
Both singlets are physical basis channels at the finite resolution.  They may
mix through the Hamiltonian when all exact quantum numbers agree, but their
identity remains visible in the state, operator, tensor-network, and
provenance records.

\begin{requirement}[Complete color multiplicity]
Every H2 basis and every H2 state export shall retain both singlet channels in
\cref{eq:two-singlets}.  A construction formed from a color-singlet \(qqq\)
state times a free gluon shall fail before Hamiltonian assembly.
\end{requirement}

The total color generator is
\begin{equation}
 \cQ^A
 =
 \sum_{i=1}^{3}t_i^A+F_g^A,
 \qquad
 (F^A)_{bc}=-if^{Abc},
\end{equation}
and every physical color tensor \(C_{\rho}\) satisfies
\begin{equation}
 \cQ^A C_{\rho}=0
 \qquad\forall A=1,\ldots,8.
\end{equation}
Generator annihilation, orthonormality, and recoupling unitarity are exact
basis gates, not posterior diagnostics.

\subsection{Longitudinal, transverse, helicity, and OAM support}

At fixed longitudinal resolution \(K\), the positive mode numbers satisfy
\begin{equation}
 \sum_{i=1}^{3}k_i+k_g=K,
 \qquad
 x_i=\frac{k_i}{K},
 \quad
 x_g=\frac{k_g}{K}>0.
\end{equation}
The intrinsic transverse momenta obey
\begin{equation}
 \sum_{i=1}^{3}\bm\kappa_{iT}+\bm\kappa_{gT}=0.
\end{equation}
The exact kinematic projection is
\begin{equation}
 \Jz
 =
 \sum_{i=1}^{3}\lambda_i+\lambda_g+
 \sum_{i=1}^{3}m_i+m_g,
 \label{eq:jz-ledger}
\end{equation}
with convention-specific oscillator and intrinsic-OAM adapters stored in the
basis identity.

The first H2 tower should support all \(L_z\) blocks reachable at its declared
\((K,N_{\max})\).  Absence of \(L_z=\pm2\) at a small benchmark resolution is
a truncation statement, not a prediction that rank-two observables vanish.

\subsection{Physical projector}

The physical finite sector is
\begin{equation}
 \Hil_{4g,\phys}^{(r)}
 =
 \cP_{\mathrm{singlet}}
 \cP_{\mathrm{antisym}}
 \cP_{\Jz,I,Q,B}
 \cP_{\COM}
 \Hil_{4g,\mathrm{kin}}^{(r)}.
\end{equation}
The projectors must be mutually compatible on the retained basis.  Their
commutator residuals
\begin{equation}
 \delta_{AB}^{(r)}
 =
 \norm{[\cP_A,\cP_B]}
\end{equation}
are exact zero when the corresponding structures should commute and explicit
truncation diagnostics otherwise.

\section{Resolution-indexed coupled Hamiltonian}

\subsection{Block mass operator}

The H2 mass-squared operator is
\begin{equation}
 \cM_{\mathrm{H2},r}^{2}
 =
 \begin{pmatrix}
  \cM_{33,r}^{2} & V_{3\leftarrow4g,r}\\
  V_{4g\leftarrow3,r} & \cM_{4g4g,r}^{2}
 \end{pmatrix},
 \qquad
 V_{3\leftarrow4g,r}=V_{4g\leftarrow3,r}^{\dagger}.
 \label{eq:block-h2}
\end{equation}
The valence block \(\cM_{33,r}^{2}\) descends from an H1 branch after removing
any effective interaction that is declared to represent explicit one-gluon
dynamics.  The \(qqqg\) block contains free invariant mass, retained induced
confinement, sector counterterms, diagonal interactions supported at the
truncation, and an explicit discrepancy basis for omitted sectors and
operators.

A useful decomposition is
\begin{align}
 \cM_{33,r}^{2}
 &={}
 \cM_{0,3,r}^{2}
 +V_{\mathrm{conf},3,r}^{\ind}
 +V_{\mathrm{spin},3,r}^{\mathrm{nonoverlap}}
 +\delta\cM_{3,r}^{2}
 +\Delta_{3,r}^{2},
 \\
 \cM_{4g4g,r}^{2}
 &={}
 \cM_{0,4g,r}^{2}
 +V_{\mathrm{conf},4g,r}^{\ind}
 +V_{\mathrm{diag},4g,r}^{\mathrm{can}}
 +\delta\cM_{4g,r}^{2}
 +\Delta_{4g,r}^{2}.
\end{align}
The label ``nonoverlap'' means that the retained valence interaction has passed
an explicit comparison against the Feshbach image of the \(qqqg\) sector.  A
legacy effective color-spin term that was declared
\texttt{ALTERNATIVE\_TO\_EXPLICIT\_QQQG\_DYNAMICS} cannot remain active by
default.

\subsection{Free invariant mass}

For either sector,
\begin{equation}
 \cM_{0,s}^{2}
 =
 \sum_{i\in s}
 \frac{\bm\kappa_{iT}^{2}+m_{i,s}^{2}}{x_i},
 \qquad
 \sum_{i\in s}x_i=1,
 \quad
 \sum_{i\in s}\bm\kappa_{iT}=0.
\end{equation}
The bare masses may be sector dependent along the finite renormalization
trajectory.  This dependence is a bookkeeping response to Fock truncation; it
is not a statement that a physical quark has different masses in different
sectors.

\subsection{Propagating quark--gluon vertex}

In light-front gauge, the canonical emission block is generated from the
quark--gluon interaction and has the schematic matrix element
\begin{equation}
 \langle\alpha_{4g}|V_{4g\leftarrow3,r}|\alpha_3\rangle
 =
 g_{34,r}^{\bare}
 \sum_{j=1}^{3}
 \cC_{j}^{a,\rho}
 \cS_{j}^{\lambda_g}
 \cL_{j}^{(r)}
 \cT_{j}^{(r)}
 \cR_{j}^{(r)}
 \cP_{j}^{\mathrm{ferm}},
 \label{eq:qg-vertex}
\end{equation}
where the factors denote color, helicity/spinor, longitudinal, transverse,
regulator, and fermionic-sign data.  The factorization in
\cref{eq:qg-vertex} is an implementation grammar, not a claim that the final
matrix element is separable after all sums and recouplings.

The vertex must preserve
\begin{equation}
 K,
 \qquad
 \Jz,
 \qquad
 I,
 \qquad
 Q,
 \qquad
 B,
\end{equation}
and its absorption block is generated by the Hilbert-space adjoint under the
same basis normalization and regulator identity.

\subsection{Diagonal one-gluon interactions}

At the one-dynamical-gluon truncation, the diagonal block can receive:
\begin{itemize}
 \item kinetic and mass terms;
 \item quark and gluon self-energy counterterms;
 \item induced interactions generated by excluded sectors;
 \item instantaneous interactions that act without changing the retained
 Fock number;
 \item regulator-dependent confinement or long-distance terms admitted by the
 selected H2 branch;
 \item current- and Ward-required contact terms.
\end{itemize}
The canonical three-gluon vertex changes the number or routing of dynamical
gluons and generally connects to sectors not fully represented by
\(qqq\oplus qqqg\).  Its omitted effects must be represented by named induced
operators and discrepancy terms, not by claiming that non-Abelian dynamics is
complete at H2.

\begin{remark}
A one-gluon Fock sector can carry adjoint color and nontrivial gluon helicity,
but it cannot by itself realize every matrix element of the nonlinear term
\(g f^{abc}A_b^iA_c^j\) in the field strength.  Operator readiness must state
which linear and nonlinear pieces are represented and which require
\(qqqgg\) or induced operators.
\end{remark}

\section{Instantaneous interactions and constrained fields}

\subsection{Origin in light-front quantization}

Light-front quantization separates dynamical and constrained field components.
Eliminating constrained components produces instantaneous fermion and gluon
interactions in the Hamiltonian.  Their exact form depends on gauge,
boundary, zero-mode, and inverse-derivative prescriptions
\cite{BrodskyPauliPinsky1998}.  In a truncated theory they are not optional
corrections: they participate in cancellations required by current
conservation and gauge consistency.

The generic inverse derivative is represented as a typed operator
\begin{equation}
 \left(\frac{1}{\partial^+}\right)_{\cB,\eps_0},
\end{equation}
where \(\cB\) is the boundary prescription and \(\eps_0\) is the endpoint or
zero-mode regulator.  Two operators with different prescriptions cannot be
identified by equal numerical arrays on one finite grid.

\subsection{Instantaneous-fermion partner}

The instantaneous-fermion term generated by the constrained spinor component
has the schematic form
\begin{equation}
 V_{\mathrm{inst},\psi}
 \sim
 g^2\,
 \bar\psi\,\gamma^i A_i
 \frac{\gamma^+}{2i\partial^+}
 \gamma^j A_j\,\psi,
 \label{eq:inst-fermion}
\end{equation}
with color, ordering, regulator, and normal-ordering data retained.  In the H2
space, some matrix elements of \cref{eq:inst-fermion} are explicit and others
represent excluded higher-gluon sectors.  The implementation shall distinguish
those cases.

\subsection{Instantaneous-gluon partner}

The constrained gauge component induces a current--current interaction of the
schematic form
\begin{equation}
 V_{\mathrm{inst},A}
 \sim
 -\frac{g^2}{2}
 J^{+a}
 \frac{1}{(\partial^+)^2}
 J^{+a},
 \label{eq:inst-gluon}
\end{equation}
where \(J^{+a}\) includes all retained quark and gluon color currents supported
by the truncation.  The same inverse-derivative and endpoint convention must
be used in the Hamiltonian, current, Ward benchmark, and any induced-operator
comparison.

\begin{requirement}[Instantaneous-partner completeness]
For every finite-order Ward benchmark, the manifest shall enumerate the
propagating and instantaneous diagrams or matrix blocks required at that
order.  Removing any required partner shall produce a nonzero signed defect.
\end{requirement}

\subsection{Zero modes and endpoint policy}

The first H2 calculation may exclude explicit gluon zero modes, but it must
record:
\begin{itemize}
 \item which modes are excluded;
 \item which inverse-derivative prescription replaces them;
 \item which symmetry or Ward residuals are sensitive to the exclusion;
 \item which induced operators or discrepancy terms are intended to represent
 their effect;
 \item the regulator trajectory along which the policy is tested.
\end{itemize}
``Zero mode absent from the basis'' is not equivalent to ``zero-mode physics is
zero.''

\section{Sector-dependent renormalization}

\subsection{Renormalization datum}

At each resolution \(r\), define
\begin{equation}
 \mathfrak R_{r}^{\mathrm{H2}}
 =
 \left(
 \cR_r,
 \theta_{3,r}^{\bare},
 \theta_{4g,r}^{\bare},
 \delta\theta_{3,r},
 \delta\theta_{4g,r},
 \delta g_{34,r},
 \{\cC_i\},
 \cS_r,
 \cZ_r,
 \Delta_r
 \right).
 \label{eq:ren-datum}
\end{equation}
The data include at least:
\begin{itemize}
 \item sector-dependent quark and gluon mass parameters;
 \item wave-function or residue renormalization data;
 \item the \(qqq\leftrightarrow qqqg\) vertex counterterm;
 \item current counterterms and induced current operators;
 \item instantaneous-interaction coefficients when the regulator splits their
 canonical relation;
 \item induced confinement and nonoverlap interactions;
 \item a truncation-discrepancy basis;
 \item the common physical renormalization conditions.
\end{itemize}
Sector-dependent renormalization is a controlled response to the asymmetric
self-energy content of a Fock truncation
\cite{Karmanov2008,Chabysheva2017}.  It is not a license to fit each sector or
observable independently.

\subsection{Calibration conditions and holdouts}

The first H2 tower may use the shared benchmark conditions
\begin{equation}
 M_N^2=M_{\mathrm{ref}}^2,
 \qquad
 F_1^p(0)=1,
 \qquad
 F_1^n(0)=0,
 \label{eq:h2-calibration}
\end{equation}
plus one declared vertex condition and one Abelianized Ward condition.
Examples of required holdouts are:
\begin{itemize}
 \item a second vertex kinematic point;
 \item a nonzero-transfer proton current;
 \item a nonzero-transfer neutron current;
 \item a second current component;
 \item a gluon momentum or helicity moment;
 \item a rotational or multiplet diagnostic;
 \item the next resolution or Fock-sector point.
\end{itemize}
The number of independent fitted conditions shall remain smaller than the
number of independent closure and prediction tests.

\subsection{Flow across the truncation tower}

Let \(r\preceq r'\) denote an enlarged basis or sector resolution.  Bare and
induced parameters flow,
\begin{equation}
 \theta_r\longmapsto\theta_{r'},
\end{equation}
while matched observables satisfy
\begin{equation}
 \delta_i(r',r)
 =
 \norm{
 R_{r'\to r}[\cO_i^{(r')}]
 -\cO_i^{(r)}
 }
 \longrightarrow0.
 \label{eq:observable-flow}
\end{equation}
The flow report separates:
\begin{enumerate}
 \item parameter drift;
 \item state-comparison residual;
 \item operator-matching residual;
 \item finite-basis Ward residual;
 \item tensor-network residual;
 \item numerical quadrature residual.
\end{enumerate}
No one category may be used to hide another.

\subsection{Naturalness and identifiability}

The renormalization trajectory shall report dimensionless naturalness
combinations, Jacobian rank, Hessian or Fisher spectrum, and profile
likelihoods for the fitted conditions.  A nearly null direction indicates that
the selected data do not distinguish two operator combinations.  The response
is to redesign the condition set or retain a correlated prior, not to assign
independent TMD parameters downstream.

\section{Ward--Takahashi and gauge-consistency program}

\subsection{Finite-basis Abelianized identity}

The first executable gauge benchmark is the Abelianized identity
\begin{equation}
 q_\mu\Gamma_r^\mu(p+q,p)
 =
 S_r^{-1}(p+q)-S_r^{-1}(p)
 +\delta_{\WTI,r}(p,q),
 \label{eq:wti}
\end{equation}
where the same Hamiltonian identity determines the self-energy, vertex,
instantaneous terms, current, and counterterms.  Light-front constructions of
currents satisfying a Ward--Takahashi identity show explicitly why
instantaneous and pair terms must accompany the reduced interaction
\cite{Marinho2008}.

The defect is decomposed as
\begin{equation}
 \delta_{\WTI,r}
 =
 \delta_{\mathrm{basis}}
 +\delta_{\mathrm{Fock}}
 +\delta_{\mathrm{endpoint}}
 +\delta_{\mathrm{zero}}
 +\delta_{\mathrm{ct}}
 +\delta_{\num}.
 \label{eq:wti-decomp}
\end{equation}
A small total defect without this decomposition is not an adequate validation
record because large terms may cancel accidentally or through overfitting.

\subsection{Current construction}

The Hamiltonian-consistent current is
\begin{equation}
 J_r^\mu
 =
 J_{(1),r}^\mu
 +J_{qqqg,r}^\mu
 +J_{\mathrm{inst},r}^\mu
 +J_{\ind,r}^\mu
 +\delta J_{\ct,r}^\mu.
 \label{eq:h2-current}
\end{equation}
Every term carries source and target sectors, momentum fibers, regulator,
current component, and Hamiltonian ancestry.  The current cannot import a
normalization from a different Hamiltonian branch or acquire a separate fit
factor at every transfer.

\subsection{From Ward--Takahashi to Slavnov--Taylor}

The non-Abelian theory requires Slavnov--Taylor identities involving ghost,
three-gluon, and composite-operator structures.  The one-gluon H2 truncation
does not contain every required sector.  Therefore the strongest allowed
status is
\begin{center}
\texttt{ABELIANIZED\_FINITE\_BASIS\_WARD\_BENCHMARKED},
\end{center}
not \texttt{FULL\_QCD\_GAUGE\_CLOSURE}.

A future non-Abelian program must add:
\begin{itemize}
 \item the necessary two-gluon and ghost-equivalent operator content or a
 rigorously matched gauge-fixed reduction;
 \item three-gluon vertex and gluon self-energy identities;
 \item BRST-compatible regulator and counterterm information;
 \item composite field-strength operator renormalization;
 \item closure under the Wilson-line and soft-sector operators.
\end{itemize}

\begin{principle}[Status monotonicity]
Passing a restricted Ward benchmark may increase only the readiness status
whose hypotheses were tested.  It cannot automatically upgrade the state,
current, gluon operator, Wilson line, or matching layer to full QCD readiness.
\end{principle}

\section{Explicit and induced descriptions}

\subsection{Feshbach reduction}

With \(P=\Pthree\) and \(Q=\Pfour\), the exact energy-dependent reduction is
\begin{equation}
 H_{\eff}(E)
 =
 PHP
 +PHQ\frac{1}{E-QHQ}QHP.
 \label{eq:feshbach-h2}
\end{equation}
The corresponding effective operator is
\begin{equation}
 \cO_{\eff}(E',E)
 =
 P\left[1+\omega^\dagger(E')\right]
 \cO
 \left[1+\omega(E)\right]P,
 \qquad
 \omega(E)=\frac{1}{E-QHQ}QHP.
 \label{eq:feshbach-op}
\end{equation}
The effective valence interaction and current are therefore linked.  Replacing
\cref{eq:feshbach-op} by \(P\cO P\) while retaining
\cref{eq:feshbach-h2} is inconsistent except in a tested limit.

\subsection{Comparison with the H1 effective branch}

The H1 color-spin interaction is compared with the Feshbach image
\begin{equation}
 V_{\mathrm{H1}}^{\eff}
 \stackrel{?}{=}
 PHQ\frac{1}{E-QHQ}QHP
 +\Delta V_{\mathrm{rem}}(E).
 \label{eq:h1-h2-match}
\end{equation}
The remainder is an explicit object.  Similar eigenvalues at one resolution do
not imply operator equivalence.  The comparison includes matrix elements,
state overlaps, current matrix elements, and withheld observables.

\subsection{Provenance relation}

The normalized provenance relation is
\begin{equation}
 \boxed{
 \text{explicit }qqqg
 \quad\Longleftrightarrow\quad
 \text{induced valence operator}
 +\text{declared remainder}
 }
 \label{eq:prov-explicit-induced}
\end{equation}
Selecting both sides as additive central mechanisms fails before numerical
evaluation.

\section{Coupled-sector tensor-network state}

\subsection{Direct-sum architecture}

The state network is
\begin{equation}
 \cN_{\mathrm{state}}
 =
 \cN_{qqq}\oplus\cN_{qqqg},
 \label{eq:sector-tn}
\end{equation}
with an explicit Fock-sector edge.  The \(qqqg\) tree may be organized as
\begin{equation}
 \left[
 (q_1\otimes q_2)_{R_{12},S_{12},L_{12},\alpha}
 \otimes q_3
 \right]_{R_{123}=8_{\rho,\lambda}}
 \otimes g_{8,\lambda_g,L_g}
 \longrightarrow
 (1,I,\Jz).
 \label{eq:qqqg-tree}
\end{equation}
The outer multiplicity \(\rho\) or \(\lambda\) is an explicit virtual label,
not an anonymous bond degeneracy.

\subsection{Hamiltonian operator network}

The block Hamiltonian is represented as an operator network with sector-
preserving and sector-changing components.  The off-diagonal vertex has one
leg in \(\cN_{qqq}\) and one in \(\cN_{qqqg}\); it may not be absorbed into a
single effective valence tensor when the explicit sector branch is selected.

Every operator tensor records:
\begin{itemize}
 \item source and target sector;
 \item color and outer multiplicity map;
 \item helicity, OAM, and \(\Jz\) selection;
 \item longitudinal and transverse basis map;
 \item regulator and counterterm identity;
 \item fermionic swap ledger;
 \item Hermitian partner identity.
\end{itemize}

\subsection{Exact tensorization and variational solution}

At the full allowed bond support, tensorization must reproduce the exact
finite-basis eigenstate.  For nested bond spaces \(\cV_{\chi_1}\subseteq
\cV_{\chi_2}\), the globally optimized energies obey
\begin{equation}
 E_{\chi_2}\leq E_{\chi_1},
 \qquad
 E_\chi\geq E_{\mathrm{exact}}.
\end{equation}
The convergence report includes
\begin{equation}
 \left
 \{
 E,
 \mathcal F_{\mathrm{state}},
 P_{4g},
 \langle x_g\rangle,
 \Delta G,
 L_q,
 L_g,
 J^\mu,
 \delta_{\WTI}
 \right\}.
 \label{eq:tn-observables}
\end{equation}
Energy-only convergence is insufficient.  A low-rank network that reproduces
the mass while erasing one color multiplicity or OAM interference fails.

\subsection{Entanglement and sector diagnostics}

The symmetry-resolved Schmidt spectrum across the Fock-sector or quark--gluon
cut is a numerical diagnostic of state complexity.  It may reveal how the
resolved gluon changes quark spin, momentum, or color correlations, as seen in
light-front Hamiltonian studies of quark--gluon entanglement
\cite{Qian2024}.  Entanglement entropy is not itself a TMD or an observable
without a declared operational map; it is a state and compression diagnostic.

\section{Eigensolvers, state tracking, and avoided crossings}

\subsection{Solver hierarchy}

Every benchmark block is solved through:
\begin{enumerate}
 \item dense Hermitian diagonalization where feasible;
 \item matrix-free Krylov or Lanczos application;
 \item the variational TTN solver;
 \item independent residual evaluation.
\end{enumerate}
The residual is
\begin{equation}
 \epsilon_{\mathrm{eig}}
 =
 \norm{\cM^2|\Psi\rangle-M^2|\Psi\rangle}.
\end{equation}
The solver certificate records orthogonality, phase convention, degeneracy
handling, and reproducibility.

\subsection{State identity across the tower}

The physical branch is tracked using exact quantum numbers, comparison maps,
subspace principal angles, current fingerprints, sector probabilities, OAM
content, and color-multiplicity content.  For a comparison map
\(I_{r\to r'}\), the overlap matrix is
\begin{equation}
 \cO_{ab}^{r'r}
 =
 \langle\Psi_a^{(r')}|
 I_{r\to r'}|\Psi_b^{(r)}\rangle.
\end{equation}
Near an avoided crossing, eigenvalue order alone is not a valid state label.
The state-tracking decision and its uncertainty are stored in the state bundle.

\subsection{Phase convention}

Each isolated state uses a deterministic phase convention, such as a selected
largest component real and positive.  Near degeneracy, a subspace frame is
fixed by parallel transport or by diagonalizing a stable reference operator.
Changing the phase convention changes amplitude coordinates but not physical
matrix elements; the adapter must nevertheless preserve member identity for
interference calculations.

\section{Microscopic quark, gluon, and current operators}

\subsection{Operator bundle}

The H2 state is accompanied by a common operator bundle
\begin{equation}
 \cB_{\mathrm{op},r}
 =
 \left
 \{
 J_r^\mu,
 T_r^{\mu\nu},
 \cO_{q,r}^{[\Gamma]},
 \cO_{g,r}^{ij},
 \cO_{qg,r},
 \cO_{\mathrm{Wilson},r}
 \right\}.
 \label{eq:operator-bundle}
\end{equation}
Every operator shares the Hamiltonian, regulator, sector, and member identity.
The bundle is versioned independently because an improved current or matching
map may be attached to the same state only through an explicit compatibility
record.

\subsection{Quark current and quark GTMD operator}

The quark bilocal acts diagonally on the active-quark number at zero skewness
in the DGLAP region.  Its H2 overlap receives contributions from both
\(qqq\) and \(qqqg\) sectors:
\begin{equation}
 W_q
 =
 W_q^{3\to3}
 +W_q^{4g\to4g}
 +\delta W_q^{\mathrm{op/match}}.
 \label{eq:quark-h2-overlap}
\end{equation}
The term \(\delta W_q^{\mathrm{op/match}}\) records constrained-field,
induced-operator, endpoint, and matching contributions not represented by the
naive diagonal overlap.

The H2 quark sector can support genuine quark--gluon correlations in twist-3
or equation-of-motion operators through \(qqq\leftrightarrow qqqg\)
interference.  BLFQ studies of twist-3 proton structure provide a concrete
example of this mechanism \cite{Zhu2024}.

\subsection{Gluon field-strength operator}

The regulated gluon operator is
\begin{equation}
 \cO_g^{ij}[U,U']
 =
 F_a^{+i}(a)
 U_{ab}[a,b]
 F_b^{+j}(b)
 U'_{ba}[b,a].
 \label{eq:gluon-op}
\end{equation}
In light-front gauge,
\begin{equation}
 F_a^{+i}
 =
 \partial^+ A_a^i
 -\partial^i A_a^+
 +g f^{abc}A_b^+A_c^i.
\end{equation}
The gauge choice simplifies but does not erase boundary links, constrained
fields, nonlinear operator terms, or renormalization.  The operator export
must declare whether it contains:
\begin{enumerate}
 \item the linear one-gluon field-strength term;
 \item nonlinear two-gluon terms;
 \item instantaneous or constrained-field terms;
 \item composite-operator counterterms;
 \item ordered Wilson links and color channel;
 \item the soft and rapidity matching status.
\end{enumerate}
At H2, the diagonal active-gluon overlap is available from the linear term.
The nonlinear term generally requires higher sectors or an induced operator.

\subsection{Transverse polarization decomposition}

The parent tensor decomposes as
\begin{equation}
 \Gamma^{ij}
 =
 \frac12\delta_T^{ij}\Gamma_{\mathrm{tr}}
 +\frac{i}{2}\epsilon_T^{ij}\Gamma_{\mathrm{hel}}
 +\Gamma_{\mathrm{ST}}^{ij},
\end{equation}
where the last term is symmetric and traceless.  The projectors are applied to
one common H2 parent and retain proton/neutron, target helicity, gluon helicity,
color multiplicity, regulator, and member identity.

\subsection{Readiness ladder}

The operator bundle distinguishes:
\begin{center}
\begin{tabular}{ll}
\toprule
Status & Meaning\\
\midrule
\texttt{FIELD\_STRENGTH\_LINEAR\_READY} & linear one-gluon matrix element validated\\
\texttt{GLUON\_HELICITY\_MATRIX\_READY} & complete finite H2 helicity parent exported\\
\texttt{REGULATED\_GTMD\_OVERLAP\_READY} & common off-forward overlap validated\\
\texttt{WILSON\_INPUT\_INTERFACE\_READY} & typed handoff exists\\
\texttt{PHYSICAL\_GLUON\_TMD\_READY} & prohibited at H2\\
\bottomrule
\end{tabular}
\end{center}

\section{Momentum, helicity, and orbital ledgers}

\subsection{Normalization and momentum}

The finite state satisfies
\begin{equation}
 P_3+P_{4g}=1,
\end{equation}
and the intrinsic plus-momentum ledger is
\begin{equation}
 \langle x_q\rangle_r+\langle x_g\rangle_r=1
 +\delta_{x,r},
 \label{eq:momentum-ledger}
\end{equation}
with \(\delta_{x,r}\) separated into basis, operator, regulator, and numerical
components.  The finite-resolution \(\langle x_g\rangle_r\) is not yet a
renormalized gluon PDF moment.

\subsection{Spin and OAM}

A canonical light-front decomposition at the model resolution is
\begin{equation}
 \frac12
 =
 \frac12\Delta\Sigma_r
 +\Delta G_r
 +L_{q,r}^{\mathrm{can}}
 +L_{g,r}^{\mathrm{can}}
 +\delta_{J,r}.
 \label{eq:spin-ledger}
\end{equation}
The operator and gauge definition is part of the identity.  The kinetic Ji
decomposition and the canonical Jaffe--Manohar decomposition are not silently
identified \cite{Ji1997,JaffeManohar1990}.  Their difference is tied to gauge
potential or torque terms and therefore to the Wilson-path structure
\cite{Hatta2012}.

The H2 state can export the finite-basis quantities in
\cref{eq:spin-ledger} as diagnostics.  A physical QCD spin decomposition
requires operator matching, scale evolution, and a declared gauge-link scheme.

\subsection{OAM ancestry}

Every OAM-sensitive GTMD or TMD projection stores the amplitude blocks whose
interference generates it.  If a projection requires \(L_z=0\) and
\(|L_z|=1\), disabling either block must remove the corresponding harmonic.
The same ancestry record is consumed by the Wilson engine so that the
absorptive phase and OAM interference cannot be disconnected.

\section{Regulated common GTMDs from the H2 state}

\subsection{Zero-skewness parent}

For a parton species \(a\), the regulated H2 parent is
\begin{equation}
 W_{a/N,\Lambda'\Lambda}^{(r)}
 (x,\kT,\DT)
 =
 \langle\Psi_N^{\Lambda';r}(P')|
 \cO_{a,r}(x,\kT,\DT)
 |\Psi_N^{\Lambda;r}(P)\rangle.
 \label{eq:h2-gtmd}
\end{equation}
The incoming and outgoing states use the same authoritative recoil map as
Volume~II.  The parent decomposes by sector and active slot,
\begin{equation}
 W_{a/N}^{(r)}
 =
 \sum_{s,t\in\{3,4g\}}
 W_{a/N}^{s\to t,(r)},
\end{equation}
with a named physical or operator source for every off-diagonal block.

At zero skewness and leading-twist DGLAP support, the canonical quark and gluon
bilocals are primarily diagonal in retained parton number.  Off-diagonal
\(3\leftrightarrow4g\) contributions belong to genuine quark--gluon,
constrained-field, induced, or matched operators and cannot be generated by
simply overlapping arrays of different dimension.

\subsection{Common reductions}

The regulated parent exposes
\begin{align}
 \mathsf T_r[W]
 &=W(x,\kT,\DT=0),\\
 \mathsf G_r[W]
 &=\int\dd^2\kT\,W(x,\kT,\DT),\\
 \mathsf P_r[W]
 &=\int\dd^2\kT\,W(x,\kT,0),\\
 \mathsf J_r[W]
 &=\int\dd x\,\dd^2\kT\,W.
\end{align}
These are regulated reductions.  Their physical identification with TMDs,
GPDs, PDFs, and local currents requires the matching layer of Volume~V.

\subsection{Gluon parent and color multiplicities}

The active-gluon parent retains both H2 color multiplicities:
\begin{equation}
 W_g
 =
 \sum_{\rho,\rho'\in\{\rho,\lambda\}}
 W_g^{\rho'\rho}.
\end{equation}
The operator may mix the multiplicities when allowed by color and all other
quantum numbers.  Collapsing the parent to a single scalar before the color and
polarization projectors would erase physical interference and invalidate the
connection to ordered-link \(f/d\) channels.

\subsection{Missing sea}

Because H2 has no explicit \(qqqq\bar q\) sector, positive-\(x\) antiquark
GTMDs are unavailable as microscopic state predictions.  The correct status is
\begin{center}
\texttt{UNAVAILABLE\_UNTIL\_H3\_PAIR\_SECTOR},
\end{center}
not a zero antiquark distribution and not an imported phenomenological curve
inside the microscopic parent.

\section{Microscopic handoff to Wilson dynamics}

\subsection{Handoff object}

The typed handoff is
\begin{equation}
 \mathfrak W_{\mathrm{handoff}}
 =
 \left(
 \begin{array}{l}
 \text{state-bundle id},\ 
 \text{Hamiltonian and coupling id},\\
 \text{regulator and resolution},\ 
 \text{sector amplitudes},\\
 \text{color multiplicities},\ 
 \text{helicity/OAM blocks},\\
 \text{field-strength operator id},\ 
 \text{spectral/cut status}
 \end{array}
 \right).
 \label{eq:wilson-handoff}
\end{equation}
The C5/C6 path, pole, cut-ledger, ordered-link, \(f/d\), antiunitary reversal,
phase-budget, and soft-overlap objects are reused unchanged.

\subsection{Cut support is not inferred from a bound state}

A discrete off-shell H2 eigenbasis has no absorptive part merely because a
small \(i\eps\) is introduced numerically.  A nonzero physical cut requires a
declared continuum or finite-volume spectral rule.  Therefore
\begin{equation}
 \text{H2 state exists}
 \not\Longrightarrow
 \text{physical cut support exists}.
\end{equation}
The highest initial status is
\texttt{MICROSCOPIC\_WILSON\_INPUT\_INTERFACE\_VALIDATED}.

\subsection{One-gluon Wilson order}

The first-order Wilson expansion is
\begin{equation}
 U_\gamma
 =1+ig\int_\gamma\dd\xi\cdot A(\xi)+\cO(g^2).
\end{equation}
H2 supplies a dynamical one-gluon state and the coupling needed to replace the
analytic C5/C6 amplitudes.  The eikonal pole orientation remains path-derived,
\begin{equation}
 \frac{1}{v\cdot\ell-i0\eta}
 =
 \PV\frac{1}{v\cdot\ell}
 +i\eta\pi\delta(v\cdot\ell).
\end{equation}
A common state and coupling must generate both Sivers-like and Boer--Mulders-
like projections; no universal scalar phase is introduced.

\subsection{Active-gluon ordered links and color channels}

For gluons, the handoff preserves the ordered pair
\begin{equation}
 [U,U']\in
 \{[+,+],[-,-],[+,-],[-,+]\}
\end{equation}
and the independent contractions
\begin{equation}
 C_f^{abc}=if^{abc},
 \qquad
 C_d^{abc}=d^{abc}.
\end{equation}
Link topology, ordering, color contraction, and process color weights remain
independent identities.  The C6 analytic \(f/d\) projectors can be applied to
the H2 parent only after the microscopic color tensor and operator ancestry are
preserved.

\subsection{Sector-completeness obligation for higher Wilson orders}

At second order,
\begin{equation}
 U_\gamma^{(2)}
 \sim
 (ig)^2\int_\gamma\int_\gamma
 \cP[A(\xi)A(\zeta)],
\end{equation}
intermediate states may require \(qqqgg\) and \(qqqq\bar q g\) sectors or
matched induced operators.  Non-Abelian exponentiation organizes connected
color structures but does not remove the need for the physical intermediate
states and counterterms \cite{Gardi2013}.  Therefore H2 cannot claim Wilson-
order convergence beyond the sector content it can represent.

\begin{requirement}[Wilson-order/Fock-order compatibility]
Every Wilson-order result shall include a manifest of the Fock sectors and
induced operators required by its intermediate states.  A missing required
sector produces \texttt{WILSON\_ORDER\_UNAVAILABLE}, not a truncated physical
prediction with an unreported omission.
\end{requirement}

\section{Soft, rapidity, and LF-to-QCD matching gate}

\subsection{Unsubtracted microscopic boundary}

The direct H2/Wilson calculation produces a regulated unsubtracted object
\begin{equation}
 W_{\mathrm{unsub}}^{\mathrm{H2}}
 (r,\cR_{\mathrm{rap}},\gamma).
\end{equation}
The physical TMD requires
\begin{equation}
 W_{\mathrm{ren}}^{\mathrm{QCD}}
 =
 Z_{\UV}
 R_{\mathrm{rap}}
 S^{-1/2}
 Z_{\mathrm{LF}\to\mathrm{QCD}}
 \otimes W_{\mathrm{unsub}}^{\mathrm{H2}}
 +\Delta_{\trunc}.
 \label{eq:h2-matching}
\end{equation}
The C6 half-soft benchmark verifies one analytic overlap cancellation.  It does
not complete \cref{eq:h2-matching}.

\subsection{Phase budget}

The imaginary part is partitioned as
\begin{equation}
 \operatorname{Im}\ln W_{\mathrm{ren}}
 =
 \operatorname{Im}\ln W_{\mathrm{unsub}}
 -\frac12\operatorname{Im}\ln S
 +\operatorname{Im}\ln Z_{\UV}
 +\operatorname{Im}\ln R_{\mathrm{rap}}
 +\operatorname{Im}\ln Z_{\mathrm{LF}\to\mathrm{QCD}}.
\end{equation}
Each component has an ownership and matching status.  An unresolved component
is not silently set to zero.

\subsection{Closed operator basis}

Evolution may begin only after the regulated operator basis closes under:
\begin{itemize}
 \item UV renormalization and operator mixing;
 \item quark/gluon singlet evolution where allowed;
 \item rapidity renormalization;
 \item soft subtraction;
 \item transverse-rank and Wilson-link classes;
 \item threshold and scheme transformations.
\end{itemize}
An H2 state with one active gluon does not by itself establish this closure.

\section{Assumption-conditioned H2 plans}

\subsection{Primary branches}

The compiler supports at least:
\begin{longtable}{L{0.18\textwidth}L{0.74\textwidth}}
\toprule
Plan & Content\\
\midrule
\endhead
\texttt{H2-PLAN-A} & Explicit \(qqqg\), resolution-refitted induced confinement,
instantaneous partners, no H1 effective color-spin term, exact/Krylov/TTN
solvers.\\
\texttt{H2-PLAN-B} & Explicit \(qqqg\), zero confinement, instantaneous
partners, no H1 effective color-spin term.\\
\texttt{H1-REFERENCE} & Read-only H1 effective color-spin branch, no explicit
\(qqqg\); used only for matched comparison.\\
\bottomrule
\end{longtable}
The branches are mutually exclusive complete theories at the declared
truncation.  They may be compared as conditional predictions but never added.

\subsection{Mandatory choices}

Every H2 assumption bundle declares:
\begin{itemize}
 \item regulator and basis tower;
 \item Fock sectors and color multiplicities;
 \item confinement branch;
 \item propagating vertex and coupling convention;
 \item instantaneous terms and zero-mode policy;
 \item counterterm basis and renormalization conditions;
 \item current and field-strength operator bundle;
 \item exact/Krylov/TTN solver route;
 \item Wilson handoff status;
 \item matching and downstream readiness.
\end{itemize}
A missing mandatory choice causes compilation failure.

\subsection{Refinement order}

The natural refinement order is
\begin{equation}
 (K,N_{\max},\cF,\chi,\text{operator order},\text{Wilson order})
 \preceq
 (K',N'_{\max},\cF',\chi',\ldots),
\end{equation}
with each axis recorded separately.  Increasing bond dimension cannot replace
adding a missing Fock sector; increasing basis size cannot replace a missing
instantaneous operator.

\section{Provenance two-complex for H2}

\subsection{Cells}

The H2 provenance complex contains:
\begin{itemize}
 \item 0-cells: sector spaces, Hamiltonians, currents, field-strength
 operators, states, Wilson objects, and results;
 \item 1-cells: emission, absorption, projection, reduction, matching,
 elimination, and handoff maps;
 \item 2-cells: explicit/induced equivalence, propagating/instantaneous Ward
 completion, count-once cut relations, and soft-overlap subtraction.
\end{itemize}

\subsection{Explicit/induced cell}

The square
\begin{equation}
\begin{tikzcd}[column sep=large,row sep=large]
 \Hil_{3}\oplus\Hil_{4g}
 \arrow[r,"\text{solve}"]
 \arrow[d,"\text{Feshbach}"']
 & {\lvert\Psi_{3+4g}\rangle}
 \arrow[d,"P+\omega"]\\
 \Hil_{3}^{\eff}
 \arrow[r,"\text{solve}"']
 & {\lvert\Psi_{3}^{\eff}\rangle}
\end{tikzcd}
\end{equation}
commutes only up to the declared remainder and operator transformation.  The
cell stores that remainder and prevents simultaneous selection of the two
presentations.

\subsection{Ward-completion cell}

A second cell records that the propagating vertex and self-energy alone do not
form the complete finite-basis current identity.  The path closes only after
instantaneous and counterterm contributions are attached.  A missing edge is
reported as a gauge-completion obstruction.

\subsection{Wilson/soft cell}

The microscopic rescattering path and the soft-subtracted path are related by
an explicit half-soft or scheme-specific subtraction cell.  The cell is
process and scheme dependent; it cannot be reused across incompatible Wilson
paths or color representations.

\subsection{Homology as an audit signal}

A nontrivial cycle may indicate:
\begin{itemize}
 \item an explicit and induced sector counted together;
 \item a propagating term missing its instantaneous partner;
 \item a Wilson cut counted twice;
 \item an unsubtracted soft overlap;
 \item a genuine alternative model branch.
\end{itemize}
It is an obstruction signal, not a physical amplitude.

\section{Convergence and uncertainty}

\subsection{Independent axes}

The H2 uncertainty and convergence manifest separates:
\begin{enumerate}
 \item parameter/calibration uncertainty;
 \item longitudinal and transverse basis truncation;
 \item Fock-sector truncation;
 \item color-multiplicity and symmetry implementation;
 \item regulator and zero-mode policy;
 \item sector-dependent renormalization;
 \item induced-interaction basis;
 \item tensor-network bond and topology truncation;
 \item eigensolver and state-tracking error;
 \item current and Ward defect;
 \item GTMD operator truncation;
 \item Wilson-order and cut-support uncertainty;
 \item soft, rapidity, and UV matching uncertainty;
 \item numerical quadrature and interpolation.
\end{enumerate}
These axes are not collapsed into one band before a common posterior is
defined.

\subsection{Correlated microscopic members}

A member label \(s\) identifies one complete state and operator bundle:
\begin{equation}
 s
 \mapsto
 \left
 \{
 H_s,
 |p\rangle_s,
 |n\rangle_s,
 J_s,
 \cO_{q,s},
 \cO_{g,s},
 \cO_{W,s}
 \right\}.
\end{equation}
Proton and neutron, quark and gluon, current and GTMD, and T-even and Wilson
projections retain the same member identity.

\subsection{Withheld validation}

At least one family from each of the following remains withheld:
\begin{itemize}
 \item nonzero-transfer current;
 \item gluon momentum or helicity moment;
 \item a second vertex kinematic point;
 \item rotational/current-component closure;
 \item next resolution or Fock sector;
 \item a GTMD reduction not used in calibration.
\end{itemize}
Failure updates the Hamiltonian, current, operator, counterterm, or discrepancy
model.  It does not create a function-specific normalization.

\section{Software architecture}

\subsection{Core objects}

\begin{longtable}{L{0.36\textwidth}L{0.56\textwidth}}
\toprule
Object & Responsibility\\
\midrule
\endhead
\texttt{H2Resolution} & Basis, regulator, endpoint, zero-mode, Fock, and
solver identity.\\
\texttt{QQQGColorMultiplicity} & Two exact singlet channels, recoupling, and
generator tests.\\
\texttt{CoupledFockBasis} & Antisymmetric \(qqq\oplus qqqg\) physical basis.\\
\texttt{CoupledHamiltonian} & Block mass operator and matrix-free action.\\
\texttt{PropagatingQGVertex} & Typed emission/absorption blocks.\\
\texttt{InstantaneousFermionTerm} & Constrained-spinor partner and inverse
\(\partial^+\) identity.\\
\texttt{InstantaneousGluonTerm} & Color-current contact interaction.\\
\RenTrajectory & Sector parameters, counterterms, conditions, flow, and
identifiability.\\
\texttt{HamiltonianCurrentBundle} & One Hamiltonian-consistent current family.\\
\WardManifest & Contributions and residual decomposition for the finite-basis
identity.\\
\texttt{CoupledSectorTTN} & Direct-sum state network with color multiplicities.\\
\texttt{H2EigensolverSuite} & Dense, Krylov, variational TTN, and tracking.\\
\texttt{GluonOperatorBundle} & Field strength, helicity matrix, rank, and
matching status.\\
\texttt{SpinMomentumLedger} & Sector, momentum, helicity, and OAM closure.\\
\WilsonHandoff & Typed microscopic adapter to C5/C6.\\
\ProvComplex & Explicit/induced, Ward, cut, and soft count-once cells.\\
\texttt{H2StateBundle} & Immutable state, operator, flow, solver, and readiness
artifact.\\
\bottomrule
\end{longtable}

\subsection{Determinism and content addressing}

Every basis, Hamiltonian, current, state, tensor network, and result receives a
content hash constructed from its scientific identity and source hashes.
Floating arrays alone do not define identity.  Regenerating a bundle from the
same inputs must reproduce ordering, phases, metadata, and manifests within the
declared deterministic contract.

\subsection{Distributed realization}

The coupled basis is naturally block sparse in exact quantum numbers, sector,
color multiplicity, and tensor-network symmetry labels.  The implementation
may distribute Hamiltonian application and contractions over blocks, but it
must preserve deterministic reductions and global residual checks.  A
parallel speedup cannot change basis ordering or state-member identity.

\section{Formal requirements}

The stable requirement identifiers for this volume are listed below.
\begin{enumerate}[label=\textbf{V9.\arabic*.}]
 \item The central H2 space is exactly \(\Hthree\oplus\Hfour\).
 \item Both \(qqqg\) color-singlet multiplicities are retained.
 \item Quark antisymmetry is imposed before operator assembly.
 \item All exact quantum numbers and regulator labels are typed.
 \item The block Hamiltonian is Hermitian with generated adjoint off-diagonal
 blocks.
 \item The H1 effective color-spin term is excluded unless overlap subtraction
 is supplied.
 \item The propagating quark--gluon vertex preserves momentum and \(\Jz\).
 \item The vertex retains emitting-quark and fermion-sign identity.
 \item Instantaneous-fermion terms use a declared inverse-derivative policy.
 \item Instantaneous-gluon terms use the same endpoint and regulator policy.
 \item Zero-mode omissions have an explicit discrepancy and validation record.
 \item Sector-dependent masses are distinguished from physical masses.
 \item A shared renormalization condition set is used across the tower.
 \item At least one vertex observable and one current family remain withheld.
 \item Hamiltonians and operators transform together under sector elimination.
 \item Explicit and induced \(qqqg\) descriptions are mutually exclusive.
 \item The finite-basis Ward identity includes all required attachments.
 \item Restricted Ward success is not labeled full Slavnov--Taylor closure.
 \item The direct-sum TTN retains Fock sector and color multiplicity.
 \item Full-bond tensorization reproduces the exact finite state.
 \item Variational energies satisfy the Rayleigh--Ritz bound.
 \item Bond convergence includes gluon, OAM, current, and Ward observables.
 \item State tracking uses overlaps and fingerprints, not eigenvalue order
 alone.
 \item The current shares the Hamiltonian identity.
 \item The gluon field-strength export states which linear/nonlinear pieces are
 present.
 \item Gluon polarization projections descend from one tensor parent.
 \item Momentum and spin ledgers report operator and matching status.
 \item Quark/gluon GTMDs descend from the same H2 state.
 \item Missing sea is reported unavailable, not zero.
 \item The Wilson handoff preserves state, coupling, color, OAM, and regulator
 identity.
 \item Physical cut support is not inferred from a discrete bound-state basis.
 \item Wilson order is gated by available Fock sectors or induced operators.
 \item Ordered gluon links and \(f/d\) contractions remain independent.
 \item Soft, rapidity, UV, and LF-to-QCD matching remain explicit gates.
 \item Assumption branches are immutable and mutually exclusive.
 \item Provenance two-cells enforce count-once semantics.
 \item Correlated proton/neutron and quark/gluon member identity is preserved.
 \item Every released result contains independent convergence axes.
 \item At least one withheld family tests the shared H2 model.
 \item No H2 result is promoted to a physical TMD without Volume V completion.
\end{enumerate}

\section{Analytic and finite-dimensional benchmarks}

\subsection{Benchmark H2-A: coupled-block Hermiticity}

Construct deterministic complex \(qqq\) and \(qqqg\) blocks and verify
\begin{equation}
 \langle\phi_4|V_{4\leftarrow3}|\psi_3\rangle
 =
 \langle V_{3\leftarrow4}\phi_4|\psi_3\rangle^*.
\end{equation}
Inject wrong conjugation, basis normalization, color multiplicity, fermion
sign, regulator, and momentum routing.

\subsection{Benchmark H2-B: color multiplicity and recoupling}

Verify exactly two \(qqqg\) singlets, total-generator annihilation,
orthonormality, recoupling unitarity, and deterministic phases.  Remove each
multiplicity in turn and demonstrate the lost state and operator channels.

\subsection{Benchmark H2-C: sector-dependent self-energy flow}

Use a two- or three-resolution coupled model in which the highest sector has a
different self-energy content.  Fit the common physical mass and show that
bare and counterterm parameters flow while the matched eigenvalue remains
stable.

\subsection{Benchmark H2-D: instantaneous Ward completion}

Evaluate \cref{eq:wti} using propagating, instantaneous, current, and
counterterm contributions.  The complete set closes to tolerance.  Removing
each required term produces a signed residual with the expected scaling.

\subsection{Benchmark H2-E: exact/Krylov/TTN agreement}

At each small tower point, compare exact diagonalization, matrix-free Krylov,
full-bond TTN, and variational fixed-bond results for energy, state overlap,
sector probability, current, OAM, and Ward residual.

\subsection{Benchmark H2-F: observable-sensitive compression}

Choose a low-rank TTN that reproduces the mass approximately but loses one
color-multiplicity interference or gluon-OAM observable.  The adaptive bond
policy must restore the required sector even if its Schmidt weight is small.

\subsection{Benchmark H2-G: spin and momentum ledgers}

Verify \cref{eq:momentum-ledger,eq:spin-ledger} at finite resolution and
separate basis, operator, matching, and numerical residuals.  Inject an
incorrect gluon momentum convention and a lost OAM block.

\subsection{Benchmark H2-H: field-strength polarization reconstruction}

Generate one common \(\Gamma^{ij}\) and reconstruct trace, antisymmetric, and
symmetric-traceless components.  Repeat in each H2 color multiplicity and for
allowed off-diagonal multiplicity interference.

\subsection{Benchmark H2-I: common-parent reduction closure}

Use the same regulated H2 GTMD parent to compute TMD, GPD, PDF, and current
routes at the forward point and nonzero transfers.  Distinguish regulated
closure from physical QCD matching.

\subsection{Benchmark H2-J: Feshbach Hamiltonian and operator equivalence}

Compare the full H2 calculation with the induced H1 description using
\cref{eq:feshbach-h2,eq:feshbach-op}.  Verify equivalence in the finite toy model
and expose the remainder in the realistic truncation.

\subsection{Benchmark H2-K: microscopic Wilson handoff}

Replace the analytic state in the C5/C6 interface with an H2 state bundle.
Verify path, pole, color, OAM, and coupling identity.  Confirm exact zero
absorption without declared cut support.

\subsection{Benchmark H2-L: Wilson-order/Fock-order gate}

Request a second-order Wilson calculation from an H2 state lacking
\(qqqgg\).  The compiler must refuse unless a valid induced operator and
matching record are supplied.

\subsection{Benchmark H2-M: assumption branches}

Compile and execute H2-PLAN-A, H2-PLAN-B, and H1-REFERENCE.  Verify distinct
identities, mutual exclusion, reproducibility, and conditional comparison
without additive mixing.

\subsection{Benchmark H2-N: state tracking through sector mixing}

Construct an avoided crossing as the \(qqq\leftrightarrow qqqg\) coupling is
varied.  Eigenvalue-order tracking must fail while overlap, current, sector,
and principal-angle tracking follow the intended branch.

\section{Mandatory negative tests}

The H2 implementation shall include stable identifiers for at least the
following failure classes:
\begin{enumerate}
 \item missing one \(qqqg\) singlet multiplicity;
 \item singlet \(qqq\) multiplied by a free gluon;
 \item nonzero total color generator;
 \item wrong adjoint-generator sign;
 \item lost quark antisymmetry;
 \item post-assembly antisymmetrization;
 \item wrong longitudinal sum;
 \item wrong transverse closure;
 \item wrong \(\Jz\);
 \item gluon zero mode admitted against policy;
 \item swapped basis or Hamiltonian scales;
 \item off-diagonal block not equal to the adjoint;
 \item wrong emitting-quark identity;
 \item wrong fermion permutation sign;
 \item wrong color multiplicity in the vertex;
 \item inconsistent endpoint regulator;
 \item missing instantaneous-fermion term;
 \item missing instantaneous-gluon term;
 \item duplicated instantaneous term;
 \item inverse-derivative policy mismatch;
 \item zero-mode omission declared physical zero;
 \item H1 effective color-spin retained with explicit \(qqqg\);
 \item Feshbach remainder silently discarded;
 \item operator not transformed with the Hamiltonian;
 \item sector-dependent bare mass interpreted as a physical mass;
 \item frozen counterterm across resolutions;
 \item separate charge normalization for each current component;
 \item fitted nonzero-transfer current used as a holdout;
 \item rank-deficient renormalization condition set hidden;
 \item propagating-only Ward claim;
 \item accidental cancellation without residual decomposition;
 \item Abelianized closure promoted to full STI closure;
 \item current imported from another Hamiltonian branch;
 \item TTN drops the Fock-sector edge;
 \item TTN merges color multiplicities;
 \item variational energy below the exact ground state;
 \item nonmonotone nested variational energies without diagnosis;
 \item mass convergence used to hide OAM failure;
 \item state tracked only by eigenvalue order;
 \item phase convention changes member identity;
 \item linear field-strength operator labeled complete;
 \item nonlinear two-gluon term silently omitted;
 \item gluon projector applied to separate ad hoc parents;
 \item quark/gluon member identities mixed;
 \item momentum sum rule uses mismatched operator conventions;
 \item canonical and kinetic OAM silently identified;
 \item absent antiquark sector reported as zero prediction;
 \item off-diagonal Fock overlap without named operator source;
 \item local recoil formula duplicates the authoritative map;
 \item Wilson handoff loses Hamiltonian coupling identity;
 \item numerical \(\eps\) treated as physical cut support;
 \item second Wilson order without sector support;
 \item ordered gluon links collapsed;
 \item \(f\) and \(d\) channels merged;
 \item process color weights assigned without process map;
 \item soft overlap omitted;
 \item soft overlap duplicated;
 \item UV matching declared zero rather than unresolved;
 \item unmatched H2 output passed to evolution;
 \item unmatched H2 output passed to a physical process;
 \item H2 output passed to nuclear composition without required bundle data;
 \item mutually exclusive assumption branches added;
 \item provenance count-once cell removed;
 \item production registry modified;
 \item accepted artifact hash changed;
 \item C5/C6 manifest changed;
 \item formalism source changed by the implementation job;
 \item hidden random seed changes a state bundle;
 \item cache reused across incompatible operator identities;
 \item distributed reduction changes deterministic ordering;
 \item one uncertainty category used to absorb another;
 \item failed holdout repaired by a TMD-specific coefficient;
 \item physical TMD readiness asserted at H2;
 \item inference readiness asserted without Volume VI gates.
\end{enumerate}

\section{Staged implementation program}

\subsection{H2.0: coupled basis and block operator}

Implement the complete \(qqq\oplus qqqg\) basis, two color multiplicities,
propagating vertex, adjoint, free and induced diagonal blocks, and exact
Hermiticity tests.

\subsection{H2.1: sector renormalization and Ward benchmark}

Add sector masses, vertex and current counterterms, instantaneous partners,
renormalization conditions, flow manifests, and the Abelianized Ward identity.

\subsection{H2.2: coupled eigensolver and tensor network}

Run exact, Krylov, and variational TTN solvers over a small tower.  Add state
tracking, bond convergence, and observable-sensitive compression tests.

\subsection{H2.3: gluon/OAM and operator exports}

Export sector probability, momentum, helicity, OAM, field-strength helicity
matrix, and regulated common-parent reductions.

\subsection{H2.4: Wilson reconnection}

Connect the microscopic state bundle to the C5/C6 engine without asserting cut
support or QCD matching.  Validate the ordered-link and \(f/d\) interfaces.

\subsection{H2.5: explicit/induced comparison}

Construct the Feshbach comparison with H1, record the remainder, and normalize
the provenance two-cell.

\subsection{H3 handoff: light sea and chiral dynamics}

The next state enlargement is
\begin{equation}
 qqq\oplus qqqg
 \longrightarrow
 qqq\oplus qqqg\oplus qqqq\bar q,
\end{equation}
with separate \(u\bar u\) and \(d\bar d\) sectors, complete color and
antisymmetry, chiral interactions, axial and pseudoscalar currents, PCAC tests,
and explicit-versus-induced sea subtraction.

\section{Risk register and safeguards}

\begin{longtable}{L{0.25\textwidth}L{0.34\textwidth}L{0.33\textwidth}}
\toprule
Risk & Failure mode & Safeguard\\
\midrule
\endhead
Finite-matrix realism & A small successful H2 block is called a QCD nucleon. &
Use readiness statuses, towers, holdouts, and explicit nonclaims.\\
Color truncation & One of the two singlets is dropped. & Exact multiplicity and
generator tests at basis construction.\\
Gauge closure by fitting & Current coefficients are tuned until the Ward defect
vanishes. & Shared Hamiltonian/current ownership and term-by-term defect
manifest.\\
Instantaneous omission & Propagating diagrams are treated as complete. &
Mandatory partner registry and ablation tests.\\
Double counting & H1 effective gluon exchange is retained with explicit
\(qqqg\). & Provenance exclusion and Feshbach comparison.\\
False sector convergence & Stable mass hides unstable gluon/OAM observables. &
Multi-observable convergence manifest.\\
Tensor-network bias & Low bond rank erases a color or OAM channel. & Required-
sector floor and observable-sensitive adaptation.\\
False cut & Numerical \(i\eps\) is interpreted as absorption. & Typed physical
cut-support gate.\\
Wilson overreach & Higher eikonal order is computed without higher Fock
support. & Wilson/Fock compatibility manifest.\\
Operator incompleteness & Linear field strength is labeled complete. & Piecewise
operator readiness and nonlinear-term gate.\\
Matching overreach & Regulated overlap is evolved as a physical TMD. & Volume V
matching gate.\\
State misidentification & An avoided crossing changes the physical branch. &
Overlap/current/subspace tracking.\\
Parameter migration & A failed gluon observable acquires its own coefficient. &
Hamiltonian/operator ownership and withheld validation.\\
Computational explosion & Full sectors become intractable. & Symmetry blocks,
TTN compression, exact small-space oracles, and distributed deterministic
application.\\
\bottomrule
\end{longtable}

\section{Final acceptance criteria}

H2 may be called a validated dynamical-gluon microscopic benchmark only when:
\begin{enumerate}
 \item the coupled basis contains complete \(qqq\) and \(qqqg\) physical
 spaces;
 \item both color singlet multiplicities and exact antisymmetry pass;
 \item the block Hamiltonian is Hermitian and reproducible;
 \item propagating and instantaneous terms share one regulator identity;
 \item the renormalization trajectory is defined at multiple resolutions;
 \item calibration and holdout conditions are frozen;
 \item the finite-basis Ward identity closes with a decomposed residual;
 \item the current is Hamiltonian consistent;
 \item exact, Krylov, and full-bond TTN states agree;
 \item variational TTN convergence includes gluon and OAM observables;
 \item state tracking survives an avoided crossing;
 \item sector, momentum, helicity, and OAM ledgers close to tolerance;
 \item the gluon field-strength helicity matrix reconstructs correctly;
 \item regulated common-parent reductions close;
 \item the explicit/induced H1/H2 comparison includes its remainder;
 \item the Wilson handoff preserves all microscopic identity;
 \item no absorptive result appears without declared cut support;
 \item higher Wilson order is refused without sector support;
 \item all matching and downstream gates remain fail-closed;
 \item all previous production artifacts and validation manifests remain
 immutable;
 \item at least one withheld current or gluon/OAM observable is successfully
 predicted from the shared H2 parameters.
\end{enumerate}

These criteria do not establish a physical nucleon or TMD.  That stronger
claim additionally requires the H3 light sea, broader Fock and basis
convergence, operator matching, QCD evolution, nuclear composition, and shared
posterior validation.

\section{Recommended implementation boundary after C9}

If C9/H2 passes the preceding gates, the next computational package should be

\begin{quote}
\textbf{C10/H3 -- fully antisymmetrized \(u\bar u\) and \(d\bar d\) sectors,
chiral dynamics, axial/PCAC currents, explicit-versus-induced sea subtraction,
and positive-\(x\) microscopic antiquark exports.}
\end{quote}

The package should preserve the H2 gluon state and add:
\begin{itemize}
 \item all three \(qqqq\bar q\) color-singlet multiplicities;
 \item complete identical-quark antisymmetry beyond a cluster basis;
 \item separate \(u\bar u\) and \(d\bar d\) sectors;
 \item pair-creation and annihilation vertices with Hermitian partners;
 \item pion/chiral effective interactions and their subtraction from explicit
 pair sectors;
 \item axial and pseudoscalar currents with a finite-basis PCAC benchmark;
 \item positive-\(x\) antiquark helicity and OAM exports;
 \item a new tensor-network sector and bond-convergence study;
 \item the first microscopic quark, antiquark, and gluon common-parent bundle.
\end{itemize}

The parallel formal volume should be Volume~X: \emph{Light-Sea Fock Sectors,
Chiral Symmetry, PCAC, and Microscopic Antiquark GTMDs}.

\section{Conclusion}

The \(qqqg\) sector is the first point at which the microscopic program contains
a gauge-field degree of freedom in the same interacting state that generates
the nucleon mass, current, color, helicity, orbital, GTMD, and Wilson inputs.
That achievement is meaningful only if the propagating sector is accompanied
by its instantaneous partners, sector-dependent counterterms, transformed
operators, Ward diagnostics, color multiplicities, and no-double-counting
relations.  Volume~IX defines those obligations as one coupled object.

The resulting H2 state is not yet continuum QCD, but it is the first state that
can replace the analytic gluon pilot by a Hamiltonian-generated amplitude.  It
also exposes the remaining physics honestly: missing light sea, incomplete
non-Abelian gauge closure, higher Wilson-order sectors, nonlinear field-
strength terms, continuum cut support, soft/rapidity completion, LF-to-QCD
matching, and nuclear composition.  When those gates are respected, the
geometric and tensor-network architecture becomes predictive in the intended
sense: assumptions define a Hamiltonian and operator bundle; the bundle
generates one correlated microscopic state; and that state supplies common
quark, gluon, current, GTMD, and Wilson projections without assigning an
independent model curve to every TMD.

\appendix

\section{Hamiltonian-term registry schema}

A Hamiltonian term record contains at least
\begin{verbatim}
{
  "term_id": "...",
  "class": "FREE | PROPAGATING | INSTANTANEOUS | INDUCED | COUNTERTERM",
  "source_sector": "qqq | qqqg",
  "target_sector": "qqq | qqqg",
  "resolution_id": "...",
  "regulator_id": "...",
  "color_map": "...",
  "outer_multiplicity_map": "...",
  "helicity_oam_rule": "...",
  "longitudinal_rule": "...",
  "transverse_rule": "...",
  "endpoint_policy": "...",
  "zero_mode_policy": "...",
  "parameter_owner": "...",
  "hermitian_partner": "...",
  "provenance_node": "...",
  "readiness": "..."
}
\end{verbatim}

\section{H2 microscopic state-bundle schema}

\begin{verbatim}
{
  "bundle_id": "...",
  "assumption_bundle_id": "...",
  "prediction_plan_id": "...",
  "hamiltonian_id": "...",
  "resolution_id": "...",
  "renormalization_trajectory_id": "...",
  "proton_neutron_member_id": "...",
  "state_tracking_id": "...",
  "sector_amplitudes": {
    "qqq": "...",
    "qqqg_rho": "...",
    "qqqg_lambda": "..."
  },
  "exact_state_hash": "...",
  "ttn_state_hash": "...",
  "bond_manifest_id": "...",
  "current_bundle_id": "...",
  "gluon_operator_bundle_id": "...",
  "ward_manifest_id": "...",
  "spin_momentum_ledger_id": "...",
  "wilson_handoff_id": "...",
  "matching_status": [
    "UV_MATCHING_REQUIRED",
    "RAPIDITY_SOFT_MATCHING_REQUIRED",
    "LF_TO_QCD_MATCHING_REQUIRED"
  ],
  "readiness": [
    "H2_COUPLED_STATE_BENCHMARKED",
    "H2_TTN_BENCHMARKED",
    "ABELIANIZED_FINITE_BASIS_WARD_BENCHMARKED"
  ]
}
\end{verbatim}

\section{Ward-closure manifest}

\begin{longtable}{L{0.24\textwidth}L{0.68\textwidth}}
\toprule
Field & Required content\\
\midrule
\endhead
Identity & Hamiltonian, current, regulator, resolution, and condition hashes.\\
Kinematics & Incoming/outgoing momentum fibers and current component.\\
Propagating pieces & Self-energy and vertex attachments.\\
Instantaneous pieces & Fermion and gluon contact terms.\\
Counterterms & Mass, wave-function, vertex, and current terms.\\
Endpoint policy & Inverse-derivative and zero-mode prescription.\\
Residuals & Total and componentwise defects with signs and tolerances.\\
Ablations & Defects after removing each required contribution.\\
Status & Restricted identity tested; forbidden stronger claims.\\
\bottomrule
\end{longtable}

\section{Microscopic Wilson-handoff schema}

\begin{verbatim}
{
  "handoff_id": "...",
  "state_bundle_id": "...",
  "hamiltonian_coupling_id": "...",
  "field_strength_operator_id": "...",
  "qqqg_color_multiplicities": ["rho", "lambda"],
  "oam_blocks": [0, -1, 1, -2, 2],
  "ordered_link_pair": "...",
  "color_contraction": "f | d | unresolved",
  "path_pole_convention_id": "...",
  "cut_support_status": "ABSENT | DECLARED_FINITE_VOLUME | CONTINUUM",
  "phase_budget_id": "...",
  "soft_overlap_status": "...",
  "wilson_order": 1,
  "fock_support_manifest": "...",
  "readiness": "MICROSCOPIC_WILSON_INPUT_INTERFACE_VALIDATED"
}
\end{verbatim}

\small
\begin{thebibliography}{99}
\raggedright

\bibitem{Dirac1949}
P.~A.~M.~Dirac,
``Forms of relativistic dynamics,''
Rev. Mod. Phys. \textbf{21}, 392 (1949).

\bibitem{BrodskyPauliPinsky1998}
S.~J.~Brodsky, H.-C.~Pauli, and S.~S.~Pinsky,
``Quantum chromodynamics and other field theories on the light cone,''
Phys. Rept. \textbf{301}, 299 (1998),
arXiv:hep-ph/9705477.

\bibitem{Vary2010}
J.~P.~Vary et al.,
``Hamiltonian light-front field theory in a basis function approach,''
Phys. Rev. C \textbf{81}, 035205 (2010),
arXiv:0905.1411.

\bibitem{Xu2025}
S.~Xu, Y.~Liu, C.~Mondal, J.~Lan, X.~Zhao, Y.~Li, and J.~P.~Vary,
``Towards a first principles light-front Hamiltonian for the nucleon,''
Phys. Lett. B \textbf{867}, 139599 (2025),
arXiv:2408.11298.

\bibitem{Mondal2025}
C.~Mondal, S.~Xu, Y.~Liu, J.~Lan, Y.~Li, X.~Zhao, and J.~P.~Vary,
``Basis light-front quantization: Advancing a first principles approach for
the nucleon,''
arXiv:2503.21372 (2025).

\bibitem{Vary2025}
J.~P.~Vary, C.~Mondal, S.~Xu, X.~Zhao, and Y.~Li,
``Nucleon structure from basis light-front quantization: Status and
prospects,''
arXiv:2512.08283 (2025).

\bibitem{Zhu2024}
Z.~Zhu et al.,
``Transverse structure of the proton beyond leading twist: A light-front
Hamiltonian approach,''
arXiv:2404.13720 (2024).

\bibitem{Qian2024}
C.~Qian, S.~Xu, Y.-G.~Yang, and X.~Zhao,
``Quark and gluon entanglement in the proton based on a light-front
Hamiltonian,''
arXiv:2412.11860 (2024).

\bibitem{Karmanov2008}
V.~A.~Karmanov, J.-F.~Mathiot, and A.~V.~Smirnov,
``Systematic renormalization scheme in light-front dynamics with Fock-space
truncation,''
Phys. Rev. D \textbf{77}, 085028 (2008),
arXiv:0801.4507.

\bibitem{Chabysheva2017}
S.~S.~Chabysheva and J.~R.~Hiller,
``Light-front \(\phi_2^4\) theory with sector-dependent mass,''
Phys. Rev. D \textbf{95}, 096016 (2017),
arXiv:1612.09331.

\bibitem{Marinho2008}
J.~A.~O.~Marinho, T.~Frederico, E.~Pace, G.~Salm\`e, and P.~U.~Sauer,
``Light-front Ward--Takahashi identity for two-fermion systems,''
Phys. Rev. D \textbf{77}, 116010 (2008),
arXiv:0805.0707.

\bibitem{Gubankova2000}
E.~Gubankova, C.-R.~Ji, and S.~R.~Cotanch,
``Flow equations for quark--gluon interactions in light-front QCD,''
Phys. Rev. D \textbf{62}, 074001 (2000),
arXiv:hep-ph/0008088.

\bibitem{BrodskyDiehlHwang2001}
S.~J.~Brodsky, M.~Diehl, and D.-S.~Hwang,
``Light-cone wavefunction representation of deeply virtual Compton
scattering,''
Nucl. Phys. B \textbf{596}, 99 (2001),
arXiv:hep-ph/0009254.

\bibitem{Meissner2009}
S.~Mei{\ss}ner, A.~Metz, and M.~Schlegel,
``Generalized parton correlation functions for a spin-1/2 hadron,''
JHEP \textbf{08}, 056 (2009),
arXiv:0906.5323.

\bibitem{Lorce2011}
C.~Lorc\'e and B.~Pasquini,
``Quark Wigner distributions and orbital angular momentum,''
Phys. Rev. D \textbf{84}, 014015 (2011),
arXiv:1106.0139.

\bibitem{Ji1997}
X.~Ji,
``Gauge-invariant decomposition of nucleon spin,''
Phys. Rev. Lett. \textbf{78}, 610 (1997),
arXiv:hep-ph/9603249.

\bibitem{JaffeManohar1990}
R.~L.~Jaffe and A.~Manohar,
``The \(g_1\) problem: Deep inelastic electron scattering and the spin of the
proton,''
Nucl. Phys. B \textbf{337}, 509 (1990).

\bibitem{Hatta2012}
Y.~Hatta,
``Notes on the orbital angular momentum of quarks in the nucleon,''
Phys. Lett. B \textbf{708}, 186 (2012),
arXiv:1111.3547.

\bibitem{Gardi2013}
E.~Gardi, J.~M.~Smillie, and C.~D.~White,
``The non-Abelian exponentiation theorem for multiple Wilson lines,''
JHEP \textbf{06}, 088 (2013),
arXiv:1304.7040.

\bibitem{Collins2011}
J.~Collins,
\emph{Foundations of Perturbative QCD},
Cambridge University Press (2011).

\end{thebibliography}

\end{document}
